\section{Article 12 - Les activités et les services}\label{article-12---les-activituxe9s-et-les-services}

\subsection{Aperçu}\label{aperuxe7u}
\subsubsection{Compétition Canadienne d’Ingénierie}\label{compuxe9tition-canadienne-dinguxe9nierie}
\begin{enumerate}
 \item L'abréviation CCI est employée pour désigner la Compétition canadienne d’ingénierie.
 \item Le nom en anglais de la CCI sera « Canadian Engineering Competition », abrégé comme CEC.
\end{enumerate}

\subsubsection{Congrès Canadien sur le Leadership en Ingénierie}\label{congruxe8s-canadien-sur-le-leadership-en-inguxe9nierie}
\begin{enumerate}
 \item L’abréviation CCLI est employée pour désigner le Congrès canadien sur le leadership en ingénierie.
 \item Le nom en anglais de la CCLI sera « Canadian Engineering Leadership Conference », abrégé comme CELC.
\end{enumerate}

\subsubsection{Congrès sur la Diversité en Ingénierie}\label{congruxe8s-sur-la-diversituxe9-en-inguxe9nierie}
\begin{enumerate}
 \item L’abréviation CDI est employée pour désigner le Congrès sur la diversité en ingénierie.
 \item Le nom en anglais de la CDI sera « Conference on Diversity in Engineering », abrégé comme CDE.
\end{enumerate}

\subsubsection{Sommet du Développement des Associations en Ingénierie}\label{sommet-du-duxe9veloppement-des-associations-en-inguxe9nierie}
\begin{enumerate}
 \item L’abréviation SDAI est employée pour désigner le Sommet du développement des associations en ingénierie.
 \item Le nom en anglais du SDAI sera “Summit on Development of Engineering Societies”, abrégé comme SDES.
\end{enumerate}

\subsubsection{Congrès sur le Développement Durable en Ingénierie}\label{congruxe8s-sur-le-duxe9veloppement-durable-en-inguxe9nierie}
\begin{enumerate}
 \item L’abréviation CDDI est employée pour désigner le Congrès sur le développement durable en ingénierie.
 \item Le nom en anglais du CDDI sera « Conference on Sustainability in Engineering », abrégé CSE.
\end{enumerate}

\subsection{L'organisation d'activités et des services}\label{lorganisation-dactivituxe9s-et-des-services}
\subsubsection{Logos}\label{logos}
\begin{enumerate}
 \item Le logo des activités et des services doit inclure, partiellement ou complètement, le logo de la Fédération, tel qu’établi par l’article 1.5.1.
\end{enumerate}

\subsubsection{Accords de Partenariat d'Activité}\label{accords-de-partenariat-dactivituxe9}
\begin{enumerate}
 \item Les spécifications des responsabilités entre la FCEG et les activités sont détaillées dans les Accords de Partenariat d'Activité de la FCEG.
\end{enumerate}

\subsubsection{Les Membres organisateurs}\label{les-membres-organisateurs}

\begin{enumerate}
 \item Seuls les membres actifs, comme décrit à l'article 4.1 de cette constitution, et le Conseil d'administration peuvent organiser les activités ou services de la Fédération.
 \item Les Membres Organisateurs sont nommés lors de l'Assemblée générale.
\end{enumerate}

\subsubsection{Les Gestionnaires d'Activités}\label{les-responsables-dactivituxe9s}

\begin{enumerate}
 \item Les Gestionnaires d’Activités (GA) sont nommés par le membre actif désigné lors de l'Assemblée générale
 \item Leur mandat s'étend de la date de leur nomination jusqu'à la résolution de tout compte recevable ou payable à la conclusion de l'activité en question
 \item Ils ont la responsabilité exclusive de tout ce qui est lié à leur activité, y compris, mais sans s'y limiter, la responsabilité financière et leur vérification par un auditeur externe, la nomination d'un comité d'organisation, la signature de documents officiels, etc.
 \item Ils rendent directement compte à leur Conseil Consultatif, au membre organisateur et au Conseil d'administration de la Fédération.
\end{enumerate}